% GENERATED FILE. Edit canonical lessons, then run npm run generate:books.
% canonical-source-hash: fnv1a64:7bcee94f81c0e69c
% canonical-lessons: GU-C303-vavajhodu, GU-C303-tratakvu, GU-C303-chapru, GU-C303-jhupdu, GU-C303-ujadvu

\chapter{The Cyclone}
\label{ch:gu-the-cyclone}
\hlchaptermodality{303}

\begin{chapteropening}
\textbf{By the end of this chapter:} \emph{I can say what a cyclone did to a village.}

Complete the last lesson of chapter 303: I can say what a cyclone did to a village.
\end{chapteropening}

\section[\texorpdfstring{\gu{વાવાઝોડું} (vāvājhoḍũ)}{વાવાઝોડું (vāvājhoḍũ)}]{\gu{વાવાઝોડું} (vāvājhoḍũ) — a cyclone}

\label{lesson:GU-C303-vavajhodu}



\begin{quote}
Before the new one: say the Gujarati for to melt, then the Gujarati for a blaze.
\end{quote}

\subsection*{You'll want to know: \gu{વાવાઝોડું}}
\textbf{\gu{વાવાઝોડું}} — \emph{vāvājhoḍũ} — \textquotedblleft{}a cyclone\textquotedblright{}.

\begin{quote}
\gu{ગઈકાલે} \gu{રાત્રે} \gu{દરિયા} \gu{કિનારે} \gu{મોટું} \gu{વાવાઝોડું} \gu{આવ્યું}.
\end{quote}

\begin{quote}
\emph{gaīkāle rātre dariyā kināre moṭũ vāvājhoḍũ āvyũ.}
\end{quote}

\begin{quote}
Last night a big cyclone hit the coast.
\end{quote}

A neuter noun.

\begin{cousinweb}
is probably built on \emph{vā}, wind, which CDIAL lists for Gujarati under Sanskrit \emph{vāyú}, \textquotedblleft{}wind\textquotedblright{} (CDIAL 11544); no dictionary we use records the rest of the word, so its history is not settled. It is neuter, and it \textbf{\gu{આવે} \gu{છે}}, comes: \textbf{\gu{વાવાઝોડું} \gu{આવ્યું}}. \textbf{\gu{તોફાન}} is any storm; a \textbf{\gu{વાવાઝોડું}} is a whirling one.
\end{cousinweb}

\begin{grammarlens}[title={What you've built}]
One more word for telling what happened.
\end{grammarlens}

\subsection*{Your turn}
\begin{itemize}
\raggedright
  \item \emph{Say it:} \emph{vāvājhoḍũ}
  \item \emph{Say it:} \emph{vāvājhoḍũ}, once more
  \item \emph{Say it:} \emph{gaīkāle rātre dariyā kināre moṭũ vāvājhoḍũ āvyũ.} — Last night a big cyclone hit the coast.
  \item \emph{Recall:} say the Gujarati for to melt, then the Gujarati for a blaze, then say \emph{vāvājhoḍũ} again
\end{itemize}

\begin{tcolorbox}[breakable,colback=teal!4,colframe=teal!35!black,title={Before you move on}]
What does \gu{વાવાઝોડું} mean? (\textquotedblleft{}A cyclone\textquotedblright{}.) Say it once more.
\end{tcolorbox}
\section[\texorpdfstring{\gu{ત્રાટકવું} (trāṭakvũ)}{ત્રાટકવું (trāṭakvũ)}]{\gu{ત્રાટકવું} (trāṭakvũ) — to strike (of a storm or an attack)}

\label{lesson:GU-C303-tratakvu}



\begin{quote}
Before the new one: say the Gujarati for a blaze, then the Gujarati for a cyclone.
\end{quote}

\subsection*{You'll want to know: \gu{ત્રાટકવું}}
\textbf{\gu{ત્રાટકવું}} — \emph{trāṭakvũ} — \textquotedblleft{}to strike (of a storm or an attack)\textquotedblright{}.

\begin{quote}
\gu{વાવાઝોડું} \gu{સવારે} \gu{ગામ} \gu{પર} \gu{ત્રાટક્યું}.
\end{quote}

\begin{quote}
\emph{vāvājhoḍũ savāre gām par trāṭakyũ.}
\end{quote}

\begin{quote}
The cyclone struck the village in the morning.
\end{quote}

Past (it …): \gu{તે} \gu{વાવાઝોડું} \gu{ત્રાટક્યું}, \emph{te vāvājhoḍũ trāṭakyũ}.

\begin{cousinweb}
\textbf{\gu{ત્રાટકવું}} (\emph{trāṭakvũ}) No dictionary we use records an older form, so its history is not settled. Storms and armies \textbf{\gu{ત્રાટકે} \gu{છે}}: they fall on a place suddenly and hard.
\end{cousinweb}

\begin{grammarlens}[title={What you've built}]
One more word for telling what happened.
\end{grammarlens}

\subsection*{Your turn}
\begin{itemize}
\raggedright
  \item \emph{Say it:} \emph{trāṭakvũ}
  \item \emph{Say it:} \emph{trāṭakvũ}, once more
  \item \emph{Say it:} \emph{vāvājhoḍũ savāre gām par trāṭakyũ.} — The cyclone struck the village in the morning.
  \item \emph{Recall:} say the Gujarati for a blaze, then the Gujarati for a cyclone, then say \emph{trāṭakvũ} again
\end{itemize}

\begin{tcolorbox}[breakable,colback=teal!4,colframe=teal!35!black,title={Before you move on}]
What does \gu{ત્રાટકવું} mean? (\textquotedblleft{}To strike (of a storm or an attack)\textquotedblright{}.) Say it once more.
\end{tcolorbox}
\section[\texorpdfstring{\gu{છાપરું} (chāprũ)}{છાપરું (chāprũ)}]{\gu{છાપરું} (chāprũ) — a thatched roof}

\label{lesson:GU-C303-chapru}



\begin{quote}
Before the new one: say the Gujarati for a cyclone, then the Gujarati for to strike.
\end{quote}

\subsection*{You'll want to know: \gu{છાપરું}}
\textbf{\gu{છાપરું}} — \emph{chāprũ} — \textquotedblleft{}a thatched roof\textquotedblright{}.

\begin{quote}
\gu{વાવાઝોડામાં} \gu{અમારું} \gu{છાપરું} \gu{ઊડી} \gu{ગયું}.
\end{quote}

\begin{quote}
\emph{vāvājhoḍāmā̃ amārũ chāprũ ūḍī gayũ.}
\end{quote}

\begin{quote}
Our roof blew off in the cyclone.
\end{quote}

A neuter noun.

\begin{cousinweb}
goes back to Sanskrit \emph{chattvará}, \textquotedblleft{}a house, a bower\textquotedblright{}, and CDIAL lists the Gujarati noun, \textquotedblleft{}a thatched roof, a cottage\textquotedblright{} (CDIAL 4976). \textbf{\gu{છત}} is a ceiling or a flat roof; a \textbf{\gu{છાપરું}} is a sloping roof of thatch or tiles.
\end{cousinweb}

\begin{grammarlens}[title={What you've built}]
One more word for telling what happened.
\end{grammarlens}

\subsection*{Your turn}
\begin{itemize}
\raggedright
  \item \emph{Say it:} \emph{chāprũ}
  \item \emph{Say it:} \emph{chāprũ}, once more
  \item \emph{Say it:} \emph{vāvājhoḍāmā̃ amārũ chāprũ ūḍī gayũ.} — Our roof blew off in the cyclone.
  \item \emph{Recall:} say the Gujarati for a cyclone, then the Gujarati for to strike, then say \emph{chāprũ} again
\end{itemize}

\begin{tcolorbox}[breakable,colback=teal!4,colframe=teal!35!black,title={Before you move on}]
What does \gu{છાપરું} mean? (\textquotedblleft{}A thatched roof\textquotedblright{}.) Say it once more.
\end{tcolorbox}
\section[\texorpdfstring{\gu{ઝૂંપડું} (jhū̃pḍũ)}{ઝૂંપડું (jhū̃pḍũ)}]{\gu{ઝૂંપડું} (jhū̃pḍũ) — a hut}

\label{lesson:GU-C303-jhupdu}



\begin{quote}
Before the new one: say the Gujarati for to strike, then the Gujarati for a thatched roof.
\end{quote}

\subsection*{You'll want to know: \gu{ઝૂંપડું}}
\textbf{\gu{ઝૂંપડું}} — \emph{jhū̃pḍũ} — \textquotedblleft{}a hut\textquotedblright{}.

\begin{quote}
\gu{ગરીબ} \gu{લોકોનાં} \gu{ઝૂંપડાં} \gu{તૂટી} \gu{ગયાં}.
\end{quote}

\begin{quote}
\emph{garīb lokonā̃ jhū̃pḍā̃ tūṭī gayā̃.}
\end{quote}

\begin{quote}
The poor people's huts were wrecked.
\end{quote}

A neuter noun.

\begin{cousinweb}
comes from \emph{*jhuppa}, a form CDIAL reconstructs, and CDIAL lists the Gujarati noun, \textquotedblleft{}a hut\textquotedblright{} (CDIAL 5403). It is neuter: \textbf{\gu{ઝૂંપડું}}, plural \textbf{\gu{ઝૂંપડાં}}.
\end{cousinweb}

\begin{grammarlens}[title={What you've built}]
One more word for telling what happened.
\end{grammarlens}

\subsection*{Your turn}
\begin{itemize}
\raggedright
  \item \emph{Say it:} \emph{jhū̃pḍũ}
  \item \emph{Say it:} \emph{jhū̃pḍũ}, once more
  \item \emph{Say it:} \emph{garīb lokonā̃ jhū̃pḍā̃ tūṭī gayā̃.} — The poor people's huts were wrecked.
  \item \emph{Recall:} say the Gujarati for to strike, then the Gujarati for a thatched roof, then say \emph{jhū̃pḍũ} again
\end{itemize}

\begin{tcolorbox}[breakable,colback=teal!4,colframe=teal!35!black,title={Before you move on}]
What does \gu{ઝૂંપડું} mean? (\textquotedblleft{}A hut\textquotedblright{}.) Say it once more.
\end{tcolorbox}
\section[\texorpdfstring{\gu{ઉજાડવું} (ujāḍvũ)}{ઉજાડવું (ujāḍvũ)}]{\gu{ઉજાડવું} (ujāḍvũ) — to lay waste}

\label{lesson:GU-C303-ujadvu}



\begin{quote}
Before the new one: say the Gujarati for a thatched roof, then the Gujarati for a hut.
\end{quote}

\subsection*{You'll want to know: \gu{ઉજાડવું}}
\textbf{\gu{ઉજાડવું}} — \emph{ujāḍvũ} — \textquotedblleft{}to lay waste\textquotedblright{}.

\begin{quote}
\gu{વાવાઝોડાએ} \gu{આખું} \gu{ગામ} \gu{ઉજાડ્યું}.
\end{quote}

\begin{quote}
\emph{vāvājhoḍāe ākhũ gām ujāḍyũ.}
\end{quote}

\begin{quote}
The cyclone laid waste to the whole village.
\end{quote}

Past, with \gu{તેણે} (he or she …): \gu{તેણે} \gu{ખેતર} \gu{ઉજાડ્યું}, \emph{teṇe khetar ujāḍyũ}.

\begin{cousinweb}
\textbf{\gu{ઉજાડવું}} (\emph{ujāḍvũ}) comes from \emph{*ujjaṭati}, a form CDIAL reconstructs, and CDIAL lists the Gujarati verb, \textquotedblleft{}to lay waste\textquotedblright{} (CDIAL 1661). A storm, a war or a flood \textbf{\gu{ઉજાડે} \gu{છે}}: it leaves a place empty.
\end{cousinweb}

\begin{grammarlens}[title={What you've built}]
One more word for telling what happened.
\end{grammarlens}

\subsection*{Your turn}
\begin{itemize}
\raggedright
  \item \emph{Say it:} \emph{ujāḍvũ}
  \item \emph{Say it:} \emph{ujāḍvũ}, once more
  \item \emph{Say it:} \emph{vāvājhoḍāe ākhũ gām ujāḍyũ.} — The cyclone laid waste to the whole village.
  \item \emph{Recall:} say the Gujarati for a thatched roof, then the Gujarati for a hut, then say \emph{ujāḍvũ} again
\end{itemize}

\begin{tcolorbox}[breakable,colback=teal!4,colframe=teal!35!black,title={Before you move on}]
What does \gu{ઉજાડવું} mean? (\textquotedblleft{}To lay waste\textquotedblright{}.) Say it once more.
\end{tcolorbox}
